\section{Expected Contributions}

\subsection{Theoretical Contribution}

The first contribution is conceptual. The project reframes the problem of human interaction with AI from trust or adoption to architecture learning. This shifts the unit of analysis from immediate reliance decisions to evolving role structures within a decision system.

\subsection{Mechanism Contribution}

The second contribution is mechanistic. Rather than claiming only that AI errors reduce trust, the project shows how AI errors may alter long-run convergence paths by generating stronger negative attraction updates than equivalent human errors. This turns algorithm aversion into a path-dependent learning problem rather than a momentary preference anomaly.

\subsection{Methodological Contribution}

The third contribution is methodological. By integrating repeated behavioral experiments with an EWA-inspired structural model, the project links experimental manipulation, behavioral trajectories, and interpretable learning parameters. In this framework, EWA is not an afterthought; it is the theoretical language of the project.

\subsection{Practical Contribution}

The fourth contribution concerns AI governance. If early AI errors shape long-run authorization regimes, then deployment strategy matters as much as average model quality. Organizations may need staged delegation, reversible autonomy, buffered pilot phases, and explicit error-repair protocols to prevent low-automation lock-in driven by early shocks.
